%% ============================================================================
%% Chapter 1: The HTTP Wire Protocol: HTTP/1.1, HTTP/2 & HTTP/3
%% ============================================================================
\chapter{The HTTP Wire Protocol: HTTP/1.1, HTTP/2 \& HTTP/3}
\label{ch:http_wire}

Every API you will ever design, consume, debug, or rescue at 3~a.m.\ ultimately reduces to bytes moving across a socket according to the rules of the Hypertext Transfer Protocol. Frameworks, gateways, and client libraries exist to hide those bytes from you---and that is precisely why production incidents so often baffle teams: when an abstraction leaks, the only ground truth left is the wire. This chapter rebuilds HTTP from the socket upward. We begin with the textual, line-oriented framing of HTTP/1.1 \cite{rfc9112}, including the exact byte layout of requests and responses, the two legal ways to delimit a message body (\texttt{Content-Length} versus \texttt{Transfer-Encoding: chunked}), persistent connections, and the head-of-line (HOL) blocking problem that made HTTP/1.1 pipelining a historical dead end. We then formalize the semantics layer \cite{rfc9110}: method properties (safe, idempotent, cacheable), the status-code taxonomy with precise selection guidance, and the essential header registry that governs content negotiation, conditional requests, caching, and cross-origin access.

From there we move to the binary era. HTTP/2 \cite{rfc9113} replaces text with a 9-byte frame header and multiplexed streams, compresses headers with HPACK \cite{rfc7541}, and introduces credit-based flow control---yet it merely relocates HOL blocking from the application layer to TCP. HTTP/3 \cite{rfc9114} completes the arc by abandoning TCP for QUIC \cite{rfc9000} over UDP, giving each stream independent delivery and folding the transport handshake into TLS~1.3 \cite{rfc8446}. Because no HTTPS request exists without TLS, we dissect the TLS~1.3 handshake message by message: \texttt{ClientHello}, \texttt{ServerHello}, \texttt{EncryptedExtensions}, \texttt{Certificate}, \texttt{CertificateVerify}, and \texttt{Finished}, along with ALPN negotiation, SNI, certificate chain validation, and the 2-RTT (TLS~1.2) versus 1-RTT (TLS~1.3) economics that drive connection setup latency.

Why does this matter in production? Because the most expensive API defects are protocol-level misunderstandings: request smuggling born of \texttt{Content-Length}/\texttt{Transfer-Encoding} ambiguity, connection churn born of a client keep-alive timeout that is one second longer than the load balancer's idle timeout, retry storms born of treating non-idempotent \texttt{POST} as safe to replay, and latency cliffs born of TCP-level HOL blocking under packet loss. An engineer who can read the wire---who can hexdump a chunked body, decode an HTTP/2 frame header by hand, and narrate a TLS~1.3 flight---debugs in minutes what others debug in days. This chapter gives you that literacy, with complete, offline, stdlib-only Python laboratories that let you observe every byte yourself.

\section{Learning Objectives \& Prerequisites}

\subsection*{Prerequisites}

This chapter assumes you have completed \textbf{Chapter~0} (Introduction), which establishes the workspace conventions: Python~3.11+, PowerShell on Windows, offline-first execution, and the fictional \emph{Northwind Robotics} universe used for all synthetic data. You should be comfortable with basic Python (functions, classes, context managers) and with running scripts from a terminal. No prior networking knowledge is assumed; sockets, TLS, and framing are built from first principles.

\subsection*{Learning Objectives}

After completing this chapter, you will be able to:

\begin{enumerate}[label=\textbf{LO1.\arabic*}, leftmargin=3.2em]
  \item Read and construct raw HTTP/1.1 messages at the byte level, including request-line/status-line grammar, CRLF framing, and both body-framing mechanisms (\texttt{Content-Length} and chunked transfer coding) \cite{rfc9112}.
  \item Classify every standard method by its \emph{safe}, \emph{idempotent}, and \emph{cacheable} properties, and select the correct status code for ambiguous production scenarios (200 vs.\ 201 vs.\ 204; 400 vs.\ 422; 401 vs.\ 403 vs.\ 404; 409 vs.\ 412; 429; 503 with \texttt{Retry-After}) \cite{rfc9110}.
  \item Explain the HTTP/2 binary framing layer bit by bit: the 9-byte frame header, all ten frame types, the stream state machine, HPACK compression, and two-level flow control \cite{rfc9113,rfc7541}.
  \item Articulate why HTTP/3 moved to QUIC over UDP, how QUIC connection IDs enable migration, and what head-of-line trade-offs remain \cite{rfc9114,rfc9000}.
  \item Narrate the TLS~1.3 handshake message by message, explain 1-RTT versus 2-RTT setup and 0-RTT resumption, and configure ALPN/SNI correctly in a Python \texttt{ssl.SSLContext} \cite{rfc8446}.
  \item Build, run, and extend stdlib-only tooling: a raw-socket HTTP/1.1 client, a chunked/keep-alive demonstration server, a TLS configuration inspector, and an HTTP/2 frame parser.
  \item Diagnose the classic wire-level failure modes---request smuggling, keep-alive race conditions, HOL blocking, and blind retries of non-idempotent methods---and design mitigations.
\end{enumerate}

\begin{keyconceptbox}
\textbf{The wire is the contract.} Specifications such as OpenAPI describe what an API \emph{should} do; the wire shows what it \emph{actually} does. Every debugging session that matters ends at the byte stream. Train yourself to always ask: \emph{what did the bytes say?}
\end{keyconceptbox}

\section{Deep Technical Mechanics \& Architectural Concepts}

\subsection{HTTP/1.1 Message Anatomy on the Wire}

HTTP/1.1 is a \emph{textual}, request--response, stateless application protocol carried over a reliable byte stream, historically TCP \cite{rfc9112}. ``Textual'' is literal: the bytes on the wire are ASCII (with body octets of arbitrary encoding), lines are terminated by the two-byte sequence CRLF (\texttt{0x0D 0x0A}), and you can read an entire exchange with nothing more than a hex dump.

\begin{definition}[HTTP/1.1 message]
An HTTP/1.1 \emph{message} is the sequence: \emph{start-line}, zero or more \emph{header fields} each terminated by CRLF, an empty line (CRLF alone) marking the end of the header section, and an optional \emph{message body} whose length is determined by the framing rules of RFC~9112~\S6. A request's start-line is the \emph{request-line} (\texttt{method SP request-target SP HTTP-version CRLF}); a response's start-line is the \emph{status-line} (\texttt{HTTP-version SP status-code SP reason-phrase CRLF}).
\end{definition}

A minimal, complete request as raw bytes---note the visible \texttt{\textbackslash r\textbackslash n} escape annotations we will use throughout for readability---looks like this:

\begin{minted}{text}
GET /robots/42 HTTP/1.1\r\n
Host: api.northwind-robotics.example\r\n
Accept: application/json\r\n
Connection: keep-alive\r\n
\r\n
\end{minted}

And the corresponding response, with a 27-byte JSON body:

\begin{minted}{text}
HTTP/1.1 200 OK\r\n
Content-Type: application/json\r\n
Content-Length: 27\r\n
Connection: keep-alive\r\n
\r\n
{"id":42,"name":"Weld-Bot"}\r\n   <- trailing CRLF shown for clarity;
                                    only the 27 body bytes are counted
\end{minted}

\begin{keyconceptbox}
\textbf{The empty line is the protocol.} The single CRLF that separates the header section from the body is the most important two bytes in HTTP/1.1. Every parser---from a browser to an embedded proxy---splits head from body at the first \texttt{CRLF CRLF} sequence. Anything ambiguous about where the body starts or ends is a security incident waiting to happen (Section~\ref{sec:pitfalls}).
\end{keyconceptbox}

\subsubsection{Message Body Framing: Content-Length vs.\ Transfer-Encoding}

Because HTTP/1.1 reuses connections, a receiver must know exactly where one message ends and the next begins. There are exactly two legal delimiting mechanisms for bodies \cite{rfc9112}:

\begin{enumerate}
  \item \textbf{\texttt{Content-Length:} $n$} --- the body is exactly $n$ octets. The receiver reads $n$ bytes and the next byte starts the next message. A message with both \texttt{Transfer-Encoding} and \texttt{Content-Length} is malformed; \texttt{Transfer-Encoding} wins, and the discrepancy is the root of \emph{request smuggling} (Section~\ref{sec:pitfalls}).
  \item \textbf{\texttt{Transfer-Encoding: chunked}} --- the body is a sequence of chunks, each prefixed by its size in hexadecimal. The exact byte layout:
\end{enumerate}

\begin{minted}{text}
HTTP/1.1 200 OK\r\n
Content-Type: text/plain\r\n
Transfer-Encoding: chunked\r\n
\r\n
7\r\n                  <- chunk size, ASCII hex (0x37 0x0D 0x0A)
Weld-Bo\r\n            <- 7 data octets + CRLF
6\r\n
t #42\r\n
0\r\n                  <- terminal zero-length chunk
\r\n                   <- end of trailers (none here)
\end{minted}

Each chunk is \texttt{chunk-size [CRLF chunk-data CRLF]}; the stream ends with a zero-size chunk followed by an optional trailer section and a final CRLF. Chunked encoding exists so a sender can stream a body whose length is not known in advance (log tails, generated CSVs, Server-Sent Events). If neither header is present in a response, the body runs until connection close---which destroys keep-alive for that connection.

\subsubsection{Connection Reuse, Keep-Alive, and Pipelining}

HTTP/1.0 closed the TCP connection after every response; HTTP/1.1 makes connections \emph{persistent by default} \cite{rfc9112}. A client declares \texttt{Connection: close} only to opt out. Reuse amortizes TCP slow start and the TLS handshake across many requests---at high request rates this is the difference between thousands and tens of requests per second per connection.

\emph{Pipelining} went further: send request~2 before response~1 arrives. It died for three compounding reasons: (1)~responses must still be delivered in request order, so one slow response blocks all later ones (application-layer HOL blocking); (2)~intermediaries (proxies, ancient load balancers) routinely mangled pipelined traffic; (3)~a non-idempotent pipelined request cannot be safely retried if the connection dies mid-pipeline. Every major browser shipped with pipelining disabled and eventually removed the code entirely. The industry answer to the same problem was HTTP/2 multiplexing.

\subsection{Method Semantics per RFC 9110}

Method semantics are not conventions; they are contract terms that intermediaries (caches, retries, pre-fetchers) rely on \cite{rfc9110}.

\begin{definition}[Safe, idempotent, cacheable]
A method is \emph{safe} if it is read-only: the client does not request, and cannot be held accountable for, any state change on the server. A method is \emph{idempotent} if repeating the identical request $n \geq 1$ times has the same \emph{intended effect on server state} as one request (side effects like logging aside). A method is \emph{cacheable} if responses to it may be stored for reuse, subject to freshness directives.
\end{definition}

\begin{table}[h]
  \centering
  \small
  \begin{tabularx}{\textwidth}{l c c c X}
    \toprule
    \textbf{Method} & \textbf{Safe} & \textbf{Idempotent} & \textbf{Cacheable} & \textbf{Intended semantics} \\
    \midrule
    GET     & yes & yes & yes$^{*}$ & Retrieve a representation of the target resource. \\
    HEAD    & yes & yes & yes       & GET without the body; metadata only. \\
    POST    & no  & no  & rare$^{\dagger}$ & Process the enclosed representation per resource semantics (create, action). \\
    PUT     & no  & yes & no        & Replace the target resource's state with the enclosed representation. \\
    PATCH   & no  & no$^{\ddagger}$ & no & Apply a partial modification described by the enclosed document. \\
    DELETE  & no  & yes & no        & Remove the target resource. \\
    OPTIONS & yes & yes & no        & Discover communication options (CORS preflight). \\
    TRACE   & yes & yes & no        & Loop-back diagnostic echo; widely disabled. \\
    CONNECT & no  & no  & no        & Establish a tunnel (e.g., TLS through a proxy). \\
    \bottomrule
  \end{tabularx}
  \caption{Method properties per RFC~9110~\S9. $^{*}$Cacheable by default heuristics. $^{\dagger}$POST responses are cacheable only with explicit freshness information. $^{\ddagger}$PATCH \emph{can} be made idempotent by construction, but the spec does not require it.}
  \label{tab:methods}
\end{table}

The engineering consequences are direct: a client library may transparently retry a \texttt{PUT} after a timeout, but must \emph{never} blindly retry a \texttt{POST}; a shared cache may serve a \texttt{GET} from cache but must invalidate on \texttt{PUT}/\texttt{PATCH}/\texttt{DELETE} to the same URI; a crawler may freely \texttt{GET} and \texttt{HEAD} but nothing else.

\subsection{Status-Code Taxonomy and Selection Guidance}

Status codes are three-digit integers grouped by class \cite{rfc9110}: \textbf{1xx} informational (interim, e.g.\ \texttt{100 Continue}), \textbf{2xx} success, \textbf{3xx} redirection, \textbf{4xx} client error (the \emph{request} is at fault), \textbf{5xx} server error (the \emph{server} failed a valid request). The discrimination between 4xx and 5xx is operationally vital: 4xx errors are caller bugs and must not page anyone; 5xx errors burn the provider's error budget.

\begin{table}[h]
  \centering
  \small
  \begin{tabularx}{\textwidth}{l l X}
    \toprule
    \textbf{Dilemma} & \textbf{Choose} & \textbf{Why} \\
    \midrule
    200 vs.\ 201 vs.\ 204 & 201 & Resource created $\Rightarrow$ 201 with \texttt{Location}; action returning content $\Rightarrow$ 200; success with nothing to say $\Rightarrow$ 204 (no body allowed). \\
    400 vs.\ 422 & 422 & 400 = malformed syntax / framing; 422 = syntactically valid, semantically wrong (schema violations). \\
    401 vs.\ 403 vs.\ 404 & varies & 401 = unauthenticated (send \texttt{WWW-Authenticate}); 403 = authenticated but forbidden; 404 may \emph{mask} existence to avoid resource enumeration. \\
    409 vs.\ 412 & varies & 409 = state conflict (duplicate, version mismatch on plain writes); 412 = a conditional header precondition (\texttt{If-Match}) failed. \\
    429 & 429 & Rate limited; always pair with \texttt{Retry-After}. \\
    503 & 503 & Transient unavailability (overload, maintenance); pair with \texttt{Retry-After}; signals ``retry later,'' unlike 500. \\
    \bottomrule
  \end{tabularx}
  \caption{Status-code selection for the scenarios production teams argue about.}
  \label{tab:statusdilemmas}
\end{table}

\subsection{Header Registry Essentials}

Headers are name--value pairs, case-insensitive, order-insensitive (with Set-Cookie the classic exception). The essentials for API engineers \cite{rfc9110}:

\begin{itemize}
  \item \textbf{\texttt{Host}} --- mandatory in HTTP/1.1 requests; enables virtual hosting and SNI-based TLS routing. Missing \texttt{Host} $\Rightarrow$ \texttt{400}.
  \item \textbf{\texttt{Accept} / \texttt{Content-Type}} --- \texttt{Accept} declares what the client can parse; \texttt{Content-Type} declares what the body \emph{is}. Both use media types with optional parameters.
  \item \textbf{Content negotiation with q-values} --- weighting in $[0,1]$: \texttt{Accept: application/json;q=1.0, text/csv;q=0.5, */*;q=0.1}. The server picks the highest-$q$ representation it can produce, or answers \texttt{406 Not Acceptable}.
  \item \textbf{Conditional requests} --- \texttt{ETag} is an opaque validator for a representation. \texttt{If-None-Match: "v17"} on a GET yields \texttt{304 Not Modified} (saving bandwidth); \texttt{If-Match: "v17"} on a PUT yields \texttt{412 Precondition Failed} on a lost update---the foundation of optimistic concurrency (Chapter~10).
  \item \textbf{\texttt{Cache-Control}} --- directives such as \texttt{no-store} (never persist), \texttt{no-cache} (revalidate every time), \texttt{max-age=60}, \texttt{private} vs.\ \texttt{public}, \texttt{must-revalidate}. \texttt{no-cache} is chronically misread as ``do not cache''; it means ``cache, but revalidate.''
  \item \textbf{CORS preflight} --- browsers send \texttt{OPTIONS} with \texttt{Origin} and \texttt{Access-Control-Request-Method}; the server answers with \texttt{Access-Control-Allow-Origin}, \texttt{-Allow-Methods}, \texttt{-Allow-Headers}, \texttt{-Max-Age}. CORS is a browser-enforced relaxation of same-origin policy, \emph{not} an authorization mechanism---non-browser clients ignore it entirely.
  \item \textbf{Connection management} --- \texttt{Connection: keep-alive|close}, \texttt{Keep-Alive: timeout=5, max=100} (hop-by-hop, stripped by proxies), and in HTTP/2 the outright prohibition of connection-specific headers.
\end{itemize}

\subsection{HTTP/2: The Binary Framing Layer}

HTTP/2 keeps the semantics of HTTP (methods, status codes, headers) but replaces the wire format with a length-prefixed binary frame protocol over a single TCP connection \cite{rfc9113}. A connection begins with the client \emph{connection preface}---the 24-octet magic string \texttt{PRI * HTTP/2.0\textbackslash r\textbackslash n\textbackslash r\textbackslash nSM\textbackslash r\textbackslash n\textbackslash r\textbackslash n} followed immediately by a \texttt{SETTINGS} frame.

\subsubsection{The 9-Byte Frame Header}

Every frame starts with a fixed 9-octet header:

\begin{minted}{text}
 +-----------------------------------------------+
 |                 Length (24)                   |
 +---------------+---------------+---------------+
 |   Type (8)    |   Flags (8)   |
 +-+-------------+---------------+-------------------------------+
 |R|                 Stream Identifier (31)                      |
 +=+=============================================================+
 |                   Frame Payload (0...)                      ...
 +---------------------------------------------------------------+
\end{minted}

That is: a \textbf{24-bit} unsigned payload length (maximum default $2^{14} = 16384$ octets, negotiated up to $2^{24}-1$ via \texttt{SETTINGS\_MAX\_FRAME\_SIZE}), an \textbf{8-bit} type, an \textbf{8-bit} flag field whose meaning is type-specific, a \textbf{1-bit reserved} field, and a \textbf{31-bit} stream identifier (0 means the frame applies to the whole connection).

\begin{table}[h]
  \centering
  \small
  \begin{tabularx}{\textwidth}{l l l X}
    \toprule
    \textbf{Type} & \textbf{Code} & \textbf{Stream} & \textbf{Purpose (key flags)} \\
    \midrule
    DATA          & \texttt{0x0} & $>$0 & Body octets (\texttt{END\_STREAM=0x1}, \texttt{PADDED=0x8}). \\
    HEADERS       & \texttt{0x1} & $>$0 & HPACK header block (\texttt{END\_STREAM=0x1}, \texttt{END\_HEADERS=0x4}, \texttt{PADDED=0x8}, \texttt{PRIORITY=0x20}). \\
    PRIORITY      & \texttt{0x2} & $>$0 & (Deprecated in RFC~9113) stream dependency/weight hint. \\
    RST\_STREAM   & \texttt{0x3} & $>$0 & Abruptly terminate one stream; 32-bit error code. \\
    SETTINGS      & \texttt{0x4} & 0 & Connection parameters; \texttt{ACK=0x1} echoes acceptance. \\
    PUSH\_PROMISE & \texttt{0x5} & $>$0 & Server push reservation; now deprecated in practice. \\
    PING          & \texttt{0x6} & 0 & 8-octet opaque liveness/RTT probe (\texttt{ACK=0x1}). \\
    GOAWAY        & \texttt{0x7} & 0 & Initiate shutdown; carries last-stream-id and error code. \\
    WINDOW\_UPDATE & \texttt{0x8} & any & Increment flow-control window by $1$..$2^{31}-1$. \\
    CONTINUATION  & \texttt{0x9} & $>$0 & Continue a header block too large for one HEADERS frame (\texttt{END\_HEADERS=0x4}). \\
    \bottomrule
  \end{tabularx}
  \caption{HTTP/2 frame types (RFC~9113~\S6).}
  \label{tab:h2frames}
\end{table}

\subsubsection{Streams, Multiplexing, and the State Machine}

A \emph{stream} is an independent, bidirectional sequence of frames within the connection, identified by a 31-bit ID: client-initiated streams use \emph{odd} IDs (1, 3, 5, \dots), server-initiated (push) streams use \emph{even} IDs. Because frames carry their stream ID, requests and responses for many streams interleave freely---one slow response no longer blocks the others. That is the multiplexing that pipelining could never deliver. Figure~\ref{fig:streamsm} gives the lifecycle.

\begin{figure}[h]
  \centering
  \begin{tikzpicture}[
      font=\small,
      state/.style={rectangle, rounded corners=2pt, draw=darkblue, fill=blue!5, minimum width=2.9cm, minimum height=0.7cm},
      trans/.style={->, >=stealth, shorten >=1pt},
      lbl/.style={font=\scriptsize\itshape, fill=white, inner sep=1pt}]
    \node[state] (idle) {idle};
    \node[state, below left=1.3cm and 1.9cm of idle] (rlocal) {reserved (local)};
    \node[state, below right=1.3cm and 1.9cm of idle] (rremote) {reserved (remote)};
    \node[state, below=3.2cm of idle] (open) {open};
    \node[state, below left=1.3cm and 1.4cm of open] (hcl) {half-closed (local)};
    \node[state, below right=1.3cm and 1.4cm of open] (hcr) {half-closed (remote)};
    \node[state, below=3.2cm of open] (closed) {closed};

    \draw[trans] (idle) -- node[lbl, right]{send HEADERS} (open);
    \draw[trans] (idle) -- node[lbl, above left]{send PUSH\_PROMISE} (rlocal);
    \draw[trans] (idle) -- node[lbl, above right]{recv PUSH\_PROMISE} (rremote);
    \draw[trans] (rlocal) -- node[lbl, below left]{send HEADERS} (hcr);
    \draw[trans] (rremote) -- node[lbl, below right]{recv HEADERS} (hcl);
    \draw[trans] (open) -- node[lbl, above left]{send END\_STREAM} (hcl);
    \draw[trans] (open) -- node[lbl, above right]{recv END\_STREAM} (hcr);
    \draw[trans] (hcl) -- node[lbl, below left]{recv END\_STREAM / RST} (closed);
    \draw[trans] (hcr) -- node[lbl, below right]{send END\_STREAM / RST} (closed);
    \draw[trans] (open) -- node[lbl, right]{RST\_STREAM either way} (closed);
    \draw[trans, bend left=20] (idle) to node[lbl, left]{RST} (closed);
    \draw[trans, bend left=20] (rlocal) to node[lbl, right]{RST} (closed);
    \draw[trans, bend right=20] (rremote) to node[lbl, left]{RST} (closed);
  \end{tikzpicture}
  \caption{HTTP/2 stream state machine (RFC~9113~\S5.1). Push-related transitions are retained for completeness; push itself is deprecated.}
  \label{fig:streamsm}
\end{figure}

\subsubsection{HPACK Header Compression}

Sending textual headers per request wastes hundreds of bytes (cookies, \texttt{User-Agent}, \texttt{Accept}) and historically defeated TCP's initial congestion window. HPACK \cite{rfc7541} compresses the header block with three mechanisms: (1)~a \textbf{static table} of 61 predefined name/value pairs (index~1 is \texttt{:authority}, index~2 is \texttt{:method: GET}, \dots); (2)~a per-connection \textbf{dynamic table} (default 4096 octets, set by \texttt{SETTINGS\_HEADER\_TABLE\_SIZE}) of recently seen headers, maintained identically on both sides so an index reference suffices on repeat; (3)~a static \textbf{Huffman code} for string literals. Because the dynamic table is ordered state, header blocks must be decoded in order---one reason header blocks can span HEADERS + CONTINUATION frames but cannot interleave with other streams' header frames. QPACK (HTTP/3) relaxes exactly this constraint.

\subsubsection{Flow Control and Push Deprecation}

HTTP/2 flow control is credit-based and hop-by-hop: each DATA octet consumes credit from \emph{both} the per-stream window and the connection-wide window (both default to 65{,}535 octets). Receivers grant more credit with \texttt{WINDOW\_UPDATE}; senders change peer stream windows with \texttt{SETTINGS\_INITIAL\_WINDOW\_SIZE}. The two-level design lets a proxy protect the whole connection from one greedy stream while still letting fast streams run. Only DATA frames count against windows. \emph{Server push}, designed to preemptively send subresources, was deprecated in RFC~9113 and disabled by browsers: in practice it wasted bandwidth pushing assets clients already had cached, and \texttt{103 Early Hints} plus good caching won.

\subsubsection{HTTP/2's Remaining Head-of-Line Problem}

Multiplexing eliminated \emph{application-layer} HOL blocking, but HTTP/2 runs over a single ordered TCP byte stream. If one TCP segment is lost, every stream stalls until retransmission---\emph{transport-layer} HOL blocking. On lossy mobile networks a 1--2\,\% packet loss rate can make HTTP/2 slower than six parallel HTTP/1.1 connections. This is the empirical motivation for HTTP/3.

\subsection{HTTP/3 and QUIC}

HTTP/3 \cite{rfc9114} keeps HTTP semantics and the frame/stream model but runs over QUIC \cite{rfc9000}, a transport built on \textbf{UDP}. Why UDP? Not because UDP is inherently faster, but because it is the only substrate that lets a userspace transport evolve: TCP is ossified in kernels and middleboxes that inspect, rate-limit, and mangle anything that is not classic TCP. QUIC reimplements reliability, congestion control, and---critically---\emph{independent} streams in userspace on top of UDP datagrams.

\begin{itemize}
  \item \textbf{Stream independence:} each QUIC stream has its own ordered delivery; a lost packet stalls only the stream(s) whose data it carried. Transport HOL blocking is gone.
  \item \textbf{Connection IDs and migration:} QUIC connections are identified by Connection IDs, not the 4-tuple. When a phone moves from Wi-Fi to LTE and its source IP changes, the connection \emph{survives}---no re-handshake, no broken downloads.
  \item \textbf{Integrated TLS~1.3:} QUIC encrypts everything above the packet number by design; the transport and crypto handshakes are fused, giving \textbf{1-RTT} connection establishment and \textbf{0-RTT} resumption for repeat clients (with replay caveats for non-idempotent requests).
  \item \textbf{QPACK:} HPACK's ordered dynamic table would reintroduce stream coupling; QPACK allows out-of-order header-block delivery with explicit acknowledgements and encoder-stream coordination.
\end{itemize}

Deployment realities: HTTP/3 requires \texttt{Alt-Svc} advertisement or HTTPS DNS records for discovery, UDP~443 must be unblocked (many enterprise networks still throttle it), CPU cost per packet is higher than kernel TCP, and every stack keeps HTTP/2+TCP as fallback. Figure~\ref{fig:versions} summarizes the three generations.

\begin{figure}[h]
  \centering
  \begin{tikzpicture}[
      font=\small,
      box/.style={rectangle, draw=darkblue, fill=blue!5, minimum width=3.6cm, align=center},
      flow/.style={->, >=stealth},
      note/.style={font=\scriptsize\itshape, align=center}]
    % HTTP/1.1 column
    \node[box] (h1title) {\textbf{HTTP/1.1} (text)};
    \node[box, below=0.35cm of h1title, minimum height=0.65cm] (h1c1) {conn 1: req$\to$resp$\to$req};
    \node[box, below=0.15cm of h1c1, minimum height=0.65cm] (h1c2) {conn 2: req$\to$resp$\to$req};
    \node[box, below=0.15cm of h1c2, minimum height=0.65cm] (h1c3) {conn 3..6 (browsers cap)};
    \node[note, below=0.25cm of h1c3, text width=3.8cm] {app-layer HOL; per-conn TCP+TLS cost};
    % HTTP/2 column
    \node[box, right=1.1cm of h1title] (h2title) {\textbf{HTTP/2} (binary)};
    \node[box, below=0.35cm of h2title, minimum height=0.65cm] (h2s1) {stream 1: HEADERS+DATA};
    \node[box, below=0.15cm of h2s1, minimum height=0.65cm] (h2s2) {stream 3: HEADERS+DATA};
    \node[box, below=0.15cm of h2s2, minimum height=0.65cm] (h2s3) {stream 5..N interleaved};
    \node[note, below=0.25cm of h2s3, text width=3.8cm] {one TCP conn; TCP-layer HOL on loss};
    % HTTP/3 column
    \node[box, right=1.1cm of h2title] (h3title) {\textbf{HTTP/3} (QUIC/UDP)};
    \node[box, below=0.35cm of h3title, minimum height=0.65cm] (h3s1) {QUIC stream 0: independent};
    \node[box, below=0.15cm of h3s1, minimum height=0.65cm] (h3s2) {QUIC stream 4: independent};
    \node[box, below=0.15cm of h3s2, minimum height=0.65cm] (h3s3) {loss stalls one stream only};
    \node[note, below=0.25cm of h3s3, text width=3.8cm] {1-RTT setup; 0-RTT resume; migration};
    % Grouping braces
    \draw[decorate, decoration={brace, amplitude=5pt}] ($(h1c1.north west)+(-0.1,0.1)$) -- ($(h1c3.south west)+(-0.1,-0.1)$) node[midway, left=6pt, note]{6 TCP conns};
    \draw[decorate, decoration={brace, amplitude=5pt}] ($(h2s1.north west)+(-0.1,0.1)$) -- ($(h2s3.south west)+(-0.1,-0.1)$) node[midway, left=6pt, note]{1 TCP conn};
    \draw[decorate, decoration={brace, amplitude=5pt}] ($(h3s1.north west)+(-0.1,0.1)$) -- ($(h3s3.south west)+(-0.1,-0.1)$) node[midway, left=6pt, note]{1 QUIC conn};
  \end{tikzpicture}
  \caption{Three generations of HTTP concurrency: parallel connections (1.1), multiplexed streams over one TCP connection (2), and independent streams over QUIC (3).}
  \label{fig:versions}
\end{figure}

\subsection{TLS 1.3: The Handshake Behind Every HTTPS Request}

HTTP semantics travel inside TLS records (record type 22 = handshake, 23 = application data). TLS~1.3 \cite{rfc8446} is a ground-up simplification of TLS~1.2: it removes RSA key transport and static DH (forward secrecy is now mandatory), removes renegotiation, shrinks the cipher list to five AEAD suites (e.g.\ \texttt{TLS\_AES\_128\_GCM\_SHA256}), and---the headline---cuts the handshake from \textbf{2-RTT} to \textbf{1-RTT}, with \textbf{0-RTT} resumption for returning clients.

\subsubsection{Message-by-Message}

Handshake message types are single octets: \texttt{ClientHello}=1, \texttt{ServerHello}=2, \texttt{EncryptedExtensions}=8, \texttt{Certificate}=11, \texttt{CertificateVerify}=15, \texttt{Finished}=20.

\begin{enumerate}
  \item \textbf{ClientHello} --- supported versions (\texttt{supported\_versions} extension listing 1.3), 32-byte random, cipher suites, \texttt{key\_share} (the client's ephemeral (EC)DHE public key, sent \emph{speculatively}---this is what buys the 1-RTT), \texttt{server\_name} (SNI), and \texttt{application\_layer\_protocol\_negotiation} (ALPN) listing e.g.\ \texttt{h2}, \texttt{http/1.1}.
  \item \textbf{ServerHello} --- chosen cipher suite, server random, and the server's \texttt{key\_share}. From this point both sides derive handshake traffic keys: \emph{everything after ServerHello is encrypted}.
  \item \textbf{EncryptedExtensions} --- responses to extensions, including the selected ALPN protocol.
  \item \textbf{Certificate} --- the server's certificate chain (leaf first, intermediates after).
  \item \textbf{CertificateVerify} --- a signature over the handshake transcript with the certificate's private key; proves possession.
  \item \textbf{Finished} --- an HMAC over the transcript using the finished key; authenticates the entire handshake. The client answers with its own \texttt{Finished}, and application data flows.
\end{enumerate}

\begin{minted}{text}
Client                                               Server

ClientHello
+ key_share, signature_algorithms,
  psk_key_exchange_modes, SNI, ALPN
                                 -------->
                                                     ServerHello
                                                     + key_share
                                            {EncryptedExtensions}
                                            {Certificate}
                                            {CertificateVerify}
                                            {Finished}
                                 <--------   ^ server flight (1 RTT total)
{Finished}
                                 -------->
[Application Data]               <------->  [Application Data]

  ^  { } = encrypted with handshake keys;  [ ] = application keys
\end{minted}

The \emph{key schedule} is a chain of HKDF extracts/expands: an \textbf{early secret} (from the PSK, or zeros) $\rightarrow$ \textbf{handshake secret} (mixed with the ECDHE shared secret) $\rightarrow$ \textbf{master secret}, from which client/server handshake traffic keys, application traffic keys, and resumption secrets are derived. Each stage is bound to the transcript hash, so tampering anywhere upstream poisons every key downstream.

\subsubsection{TLS 1.2 vs.\ 1.3: The Round-Trip Economics}

\begin{table}[h]
  \centering
  \small
  \begin{tabularx}{\textwidth}{l l X}
    \toprule
    \textbf{Property} & \textbf{TLS 1.2} & \textbf{TLS 1.3} \\
    \midrule
    Full handshake & 2-RTT & 1-RTT \\
    Resumed handshake & 1-RTT & 0-RTT (replayable early data!) \\
    Key exchange & RSA transport allowed & (EC)DHE only; forward secrecy mandatory \\
    Handshake encryption & mostly cleartext & everything after ServerHello encrypted \\
    Cipher suites & dozens, incl.\ CBC/RC4 legacy & five AEAD suites only \\
    \bottomrule
  \end{tabularx}
  \caption{TLS 1.2 versus TLS 1.3 at a glance \cite{rfc8446}.}
  \label{tab:tls}
\end{table}

At 80\,ms RTT (typical cross-continent), a cold TLS~1.2 + TCP setup costs $\approx$240\,ms before the first request byte; TLS~1.3 cuts this to $\approx$160\,ms, QUIC to $\approx$80\,ms, and 0-RTT to zero. This is why keep-alive hit rate is a first-class SLO (see the cost template in the standard addendum).

\subsubsection{SNI, ALPN, and Certificate Chain Validation}

\textbf{SNI} (\texttt{server\_name}) lets one IP host many TLS identities; the server picks the certificate matching the requested name. \textbf{ALPN} selects the application protocol inside the handshake---\texttt{h2} or \texttt{http/1.1}---so no extra round trip is needed to discover HTTP/2 support; an ALPN mismatch (client offers only \texttt{h2}, server speaks only \texttt{http/1.1}) fails the connection with \texttt{no\_application\_protocol}. \textbf{Chain validation} walks leaf $\rightarrow$ intermediates $\rightarrow$ a locally trusted root, checking at each hop: signatures, validity windows (\texttt{notBefore}/\texttt{notAfter}), hostname match against SANs, revocation where enforced, and that intermediates carry the \texttt{basicConstraints: CA=TRUE} extension. Session \textbf{resumption} (tickets issued in \texttt{NewSessionTicket} post-handshake) skips the certificate flight; in TLS~1.3 resumption \emph{is} a PSK handshake, optionally with fresh DHE.

\begin{figure}[h]
  \centering
  \begin{tikzpicture}[font=\small,
      msg/.style={->, >=stealth},
      lbl/.style={font=\scriptsize, midway, above, sloped},
      life/.style={dashed, gray}]
    \node[font=\bfseries] (c) at (0,0) {Client};
    \node[font=\bfseries] (s) at (9.5,0) {Server};
    \draw[life] (c.south) -- ++(0,-7.6);
    \draw[life] (s.south) -- ++(0,-7.6);
    \draw[msg, darkblue] (0,-0.7)  -- node[lbl]{ClientHello + key\_share + SNI + ALPN} (9.5,-0.7);
    \draw[msg, darkblue] (9.5,-1.7) -- node[lbl]{ServerHello + key\_share} (0,-1.7);
    \draw[msg, darkblue] (9.5,-2.5) -- node[lbl]{\{EncryptedExtensions\}} (0,-2.5);
    \draw[msg, darkblue] (9.5,-3.3) -- node[lbl]{\{Certificate\}} (0,-3.3);
    \draw[msg, darkblue] (9.5,-4.1) -- node[lbl]{\{CertificateVerify\}} (0,-4.1);
    \draw[msg, darkblue] (9.5,-4.9) -- node[lbl]{\{Finished\}} (0,-4.9);
    \draw[msg, packtorange] (0,-5.7) -- node[lbl]{\{Finished\}  [+ 0-RTT early data if resuming]} (9.5,-5.7);
    \draw[msg, packtorange, double] (0,-6.6) -- node[lbl]{[Application Data --- HTTP/2 or HTTP/1.1 per ALPN]} (9.5,-6.6);
    \node[note, font=\scriptsize\itshape] at (4.75,-7.35) {1-RTT total: client sends app data after one round trip};
  \end{tikzpicture}
  \caption{TLS 1.3 full handshake flight diagram (braces = encrypted under handshake keys).}
  \label{fig:tls13}
\end{figure}

\section{Production-Ready Code Implementation}

All listings are Python~3.11+, fully typed, stdlib-only, offline (localhost), and deterministic. Run them from PowerShell with \texttt{python <file>.py}; each starts its own local server, exchanges traffic, prints the exact bytes, and shuts down cleanly.

\subsection{Listing 1 --- Raw-Socket HTTP/1.1 Client with Byte-Level Tracing}

\textbf{Design decisions.} We deliberately avoid \texttt{http.client} and open a raw \texttt{socket}: the pedagogical goal is to see that an HTTP request \emph{is} bytes. The server side is the stdlib \texttt{http.server} running in a daemon thread on an ephemeral port (\texttt{("127.0.0.1", 0)} --- the OS assigns a free port, eliminating port-collision flakiness). \textbf{Defensive notes:} the escaped-byte view uses \texttt{bytes.decode("ascii", "backslashreplace")}-style rendering so arbitrary binary never corrupts the terminal; the client sets an explicit timeout so a wedged server fails fast instead of hanging; the response is read until \texttt{Content-Length} bytes of body arrive, never with a blind \texttt{recv} loop (which would block forever on a keep-alive connection).

\begin{minted}{python}
"""Raw-socket HTTP/1.1 client against a local stdlib server.

Prints the exact bytes sent and received, with CRLF framing rendered
visibly. Offline, deterministic, stdlib only (Python 3.11+).
"""
from __future__ import annotations

import socket
import threading
from http.server import BaseHTTPRequestHandler, ThreadingHTTPServer

HOST = "127.0.0.1"
BODY = b'{"id":42,"name":"Weld-Bot"}'


def visible(data: bytes) -> str:
    """Render bytes with escapes so CRLF framing is visible."""
    return "".join(
        {0x0D: r"\r", 0x0A: r"\n"}.get(b) or (chr(b) if 32 <= b < 127 else rf"\x{b:02x}")
        for b in data
    )


class RobotHandler(BaseHTTPRequestHandler):
    protocol_version = "HTTP/1.1"  # enable keep-alive responses

    def do_GET(self) -> None:  # noqa: N802 (stdlib naming)
        if self.path != "/robots/42":
            self.send_error(404, "Not Found")
            return
        self.send_response(200)
        self.send_header("Content-Type", "application/json")
        self.send_header("Content-Length", str(len(BODY)))
        self.end_headers()
        self.wfile.write(BODY)

    def log_message(self, fmt: str, *args: object) -> None:
        pass  # keep stdout deterministic


def raw_get(host: str, port: int, path: str) -> tuple[bytes, bytes]:
    request = (
        f"GET {path} HTTP/1.1\r\n"
        f"Host: {host}:{port}\r\n"
        "Accept: application/json\r\n"
        "Connection: close\r\n"  # close => read-until-EOF is unambiguous
        "\r\n"
    ).encode("ascii")
    with socket.create_connection((host, port), timeout=5.0) as sock:
        sock.sendall(request)
        chunks: list[bytes] = []
        while True:
            chunk = sock.recv(4096)
            if not chunk:
                break
            chunks.append(chunk)
    return request, b"".join(chunks)


def main() -> None:
    server = ThreadingHTTPServer((HOST, 0), RobotHandler)
    thread = threading.Thread(target=server.serve_forever, daemon=True)
    thread.start()
    try:
        port = server.server_address[1]
        sent, received = raw_get(HOST, port, "/robots/42")
        print("=== BYTES SENT ===")
        print(visible(sent))
        print("=== BYTES RECEIVED ===")
        print(visible(received))
        assert b"HTTP/1.1 200 OK\r\n" in received
        assert received.endswith(BODY)
        print("=== OK: status line and body verified ===")
    finally:
        server.shutdown()
        server.server_close()
        thread.join(timeout=5.0)


if __name__ == "__main__":
    main()
\end{minted}

\subsection{Listing 2 --- Chunked Transfer Encoding and Keep-Alive, By Hand}

\textbf{Design decisions.} Here we implement \emph{both} ends at the socket level: a tiny \texttt{socketserver} server that emits a \texttt{Transfer-Encoding: chunked} response (three chunks) and honors keep-alive for a second request on the same connection, and a client that parses chunks incrementally. \textbf{Defensive notes:} the chunk parser validates hex sizes and enforces a maximum chunk size (a classic fuzzing target); every socket has a timeout; the server handles \texttt{Connection: close} correctly; all threads are daemons and joined.

\begin{minted}{python}
"""Chunked transfer encoding + keep-alive demo over raw sockets.

The server streams a body as three chunks, then serves a second request
on the SAME connection (keep-alive). The client parses chunks by hand.
Offline, deterministic, stdlib only (Python 3.11+).
"""
from __future__ import annotations

import socket
import socketserver
import threading

HOST = "127.0.0.1"
MAX_CHUNK = 1 << 20  # 1 MiB sanity bound against hostile size lines


def recv_until(sock: socket.socket, marker: bytes,
               limit: int = 65536) -> tuple[bytes, bytes]:
    """Read until marker (inclusive); return (head, leftover bytes).

    Keeping the leftover is essential: a single recv() may return header
    and body bytes together, and discarding the tail desyncs framing.
    """
    buf = b""
    while marker not in buf:
        chunk = sock.recv(4096)
        if not chunk:
            raise ConnectionError("EOF before marker")
        buf += chunk
        if len(buf) > limit:
            raise ValueError("header section exceeds limit")
    head, _, rest = buf.partition(marker)
    return head + marker, rest


def read_chunked_body(sock: socket.socket, buffered: bytes) -> bytes:
    """Parse a chunked body. `buffered` holds bytes already read past
    the header terminator (may be empty). Returns decoded body."""
    body = bytearray()
    pending = bytearray(buffered)

    def read_line() -> bytes:
        while b"\r\n" not in pending:
            data = sock.recv(4096)
            if not data:
                raise ConnectionError("EOF inside chunked body")
            pending.extend(data)
        line, _, rest = bytes(pending).partition(b"\r\n")
        pending.clear()
        pending.extend(rest)
        return line

    while True:
        size = int(read_line().split(b";")[0], 16)  # ignore extensions
        if size > MAX_CHUNK:
            raise ValueError(f"chunk size {size} exceeds bound")
        if size == 0:
            while read_line() != b"":  # consume trailers
                pass
            return bytes(body)
        while len(pending) < size + 2:  # data + trailing CRLF
            data = sock.recv(4096)
            if not data:
                raise ConnectionError("EOF inside chunk data")
            pending.extend(data)
        body.extend(pending[:size])
        if pending[size : size + 2] != b"\r\n":
            raise ValueError("chunk missing trailing CRLF")
        del pending[: size + 2]


class ChunkedHandler(socketserver.BaseRequestHandler):
    def handle(self) -> None:
        sock: socket.socket = self.request
        sock.settimeout(5.0)
        for _ in range(2):  # serve two requests on one connection
            head, _ = recv_until(sock, b"\r\n\r\n")  # requests are head-only
            request_line = head.split(b"\r\n", 1)[0].decode("ascii")
            keep_alive = b"connection: close" not in head.lower()
            if request_line.startswith("GET /stream "):
                parts = [b"Weld-", b"Bot ", b"#42 online"]
                payload = b"".join(
                    b"%X\r\n%s\r\n" % (len(p), p) for p in parts
                ) + b"0\r\n\r\n"
                sock.sendall(
                    b"HTTP/1.1 200 OK\r\n"
                    b"Content-Type: text/plain\r\n"
                    b"Transfer-Encoding: chunked\r\n\r\n" + payload
                )
            else:
                body = b"pong"
                sock.sendall(
                    b"HTTP/1.1 200 OK\r\nContent-Length: 4\r\n\r\n" + body
                )
            if not keep_alive:
                break


def main() -> None:
    server = socketserver.ThreadingTCPServer((HOST, 0), ChunkedHandler)
    server.daemon_threads = True
    thread = threading.Thread(target=server.serve_forever, daemon=True)
    thread.start()
    try:
        port = server.server_address[1]
        with socket.create_connection((HOST, port), timeout=5.0) as sock:
            sock.sendall(
                b"GET /stream HTTP/1.1\r\nHost: localhost\r\n\r\n"
            )
            head, buffered = recv_until(sock, b"\r\n\r\n")
            print("=== RESPONSE HEAD (request 1) ===")
            print(head.decode("ascii"), end="")
            body = read_chunked_body(sock, buffered)
            print("=== DECODED CHUNKED BODY ===")
            print(repr(body))
            # Second request on the SAME socket: keep-alive in action.
            sock.sendall(
                b"GET /ping HTTP/1.1\r\nHost: localhost\r\n"
                b"Connection: close\r\n\r\n"
            )
            # Server closes after this response: read until EOF.
            reply2 = b""
            while True:
                chunk = sock.recv(4096)
                if not chunk:
                    break
                reply2 += chunk
            print("=== RESPONSE (request 2, same connection) ===")
            print(reply2.decode("ascii"))
        assert body == b"Weld-Bot #42 online"
        assert b"pong" in reply2
        print("=== OK: chunked decode + keep-alive reuse verified ===")
    finally:
        server.shutdown()
        server.server_close()
        thread.join(timeout=5.0)


if __name__ == "__main__":
    main()
\end{minted}

\subsection{Listing 3 --- TLS Inspector: Context Configuration, ALPN, and Ciphers}

\textbf{Design decisions.} Generating a self-signed certificate purely from stdlib Python is not supported (the \texttt{ssl} module consumes but does not mint certificates), so the lab (Section~\ref{sec:lab}) mints one with a documented \texttt{openssl.exe} command. This listing is split into two halves that work with or without that certificate: (1)~a \emph{certificate fixture decoder} using the documented stdlib hook \texttt{ssl.\_ssl.\_test\_decode\_cert}, which parses a PEM file without opening any connection; (2)~a \emph{context auditor} that builds a hardened \texttt{SSLContext}, pins TLS~1.3, sets ALPN to \texttt{h2}/\texttt{http/1.1}, and prints the effective cipher list --- all offline. If the PEM fixture from the lab exists next to the script, the decoded fields (subject, issuer, SANs, validity) are printed too; otherwise the script explains exactly how to create it and still runs to completion.

\begin{minted}{python}
"""TLS inspector: certificate decoding, context hardening, ALPN, ciphers.

Offline, stdlib only (Python 3.11+). If a PEM fixture (created in the
chapter lab with openssl) exists beside this script, its fields are
decoded with ssl._ssl._test_decode_cert; otherwise the context audit
still runs and the script prints the fixture instructions.
"""
from __future__ import annotations

import ssl
from pathlib import Path

CERT_PATH = Path(__file__).with_name("northwind_lab.pem")


def decode_certificate(path: Path) -> None:
    """Parse a PEM certificate file into a dict and print key fields."""
    decoded = ssl._ssl._test_decode_cert(str(path))  # stdlib test hook
    print("=== CERTIFICATE FIELDS ===")
    for field in ("subject", "issuer", "notBefore", "notAfter"):
        print(f"{field}: {decoded.get(field)}")
    san = decoded.get("subjectAltName", ())
    print(f"subjectAltName: {san}")
    serial = decoded.get("serialNumber", "<unknown>")
    print(f"serialNumber: {serial}")


def build_client_context() -> ssl.SSLContext:
    """Hardened TLS 1.3 client context with ALPN."""
    ctx = ssl.SSLContext(ssl.PROTOCOL_TLS_CLIENT)
    ctx.minimum_version = ssl.TLSVersion.TLSv1_3
    ctx.check_hostname = True
    ctx.verify_mode = ssl.CERT_REQUIRED
    ctx.set_alpn_protocols(["h2", "http/1.1"])
    if CERT_PATH.exists():
        ctx.load_verify_locations(cafile=str(CERT_PATH))
    return ctx


def audit_ciphers(ctx: ssl.SSLContext) -> None:
    """Print the negotiated-cipher candidates of a context."""
    print("=== TLS 1.3+ CIPHER SUITES AVAILABLE ===")
    for cipher in ctx.get_ciphers():
        print(
            f"{cipher['name']:<32} protocol={cipher['protocol']:<8} "
            f"bits={cipher['strength_bits']}"
        )


def main() -> None:
    if CERT_PATH.exists():
        decode_certificate(CERT_PATH)
    else:
        print(f"NOTE: {CERT_PATH.name} not found; run the chapter lab's "
              "openssl command to create it, then re-run.")
    ctx = build_client_context()
    print(f"minimum_version: {ctx.minimum_version}")
    print(f"ALPN offered: h2, http/1.1")
    audit_ciphers(ctx)
    # NOTE: OpenSSL's cipher API may omit TLS 1.3 suites (they are not
    # configurable); an empty/short list here is expected, not an error.
    assert ctx.minimum_version == ssl.TLSVersion.TLSv1_3
    print("=== OK: context hardened and cipher audit complete ===")


if __name__ == "__main__":
    main()
\end{minted}

\subsection{Listing 4 --- HTTP/2 Frame Parser over a Captured Byte Stream}

\textbf{Design decisions.} The parser consumes the 9-octet header with \texttt{struct} (network byte order), validates the 31-bit stream mask, decodes the type-specific flags for DATA/HEADERS/SETTINGS, and prints a table. The ``captured stream'' is an embedded fixture built byte-by-byte in the script itself --- a client connection preface plus SETTINGS, SETTINGS-ACK, HEADERS, and DATA frames --- so the output is fully deterministic and reproducible offline. \textbf{Defensive notes:} the parser raises \texttt{ValueError} on truncation rather than slicing garbage, enforces \texttt{SETTINGS\_MAX\_FRAME\_SIZE} ($2^{14}$ default), and treats unknown frame types as skip-able per RFC~9113~\S5.5 (extensions must be ignorable).

\begin{minted}{python}
"""HTTP/2 frame parser: decode a captured byte stream, stdlib only.

The fixture is constructed in-script (client preface + SETTINGS,
SETTINGS-ACK, HEADERS, DATA) so results are deterministic and offline.
"""
from __future__ import annotations

import struct
from dataclasses import dataclass

MAX_FRAME_SIZE = 2**14  # RFC 9113 default SETTINGS_MAX_FRAME_SIZE
CLIENT_PREFACE = b"PRI * HTTP/2.0\r\n\r\nSM\r\n\r\n"

FRAME_TYPES = {
    0x0: "DATA", 0x1: "HEADERS", 0x2: "PRIORITY", 0x3: "RST_STREAM",
    0x4: "SETTINGS", 0x5: "PUSH_PROMISE", 0x6: "PING", 0x7: "GOAWAY",
    0x8: "WINDOW_UPDATE", 0x9: "CONTINUATION",
}
FLAG_NAMES = {
    0x0: {0x1: "END_STREAM", 0x8: "PADDED"},
    0x1: {0x1: "END_STREAM", 0x4: "END_HEADERS", 0x8: "PADDED",
          0x20: "PRIORITY"},
    0x4: {0x1: "ACK"},
    0x6: {0x1: "ACK"},
    0x9: {0x4: "END_HEADERS"},
}


@dataclass(frozen=True)
class Frame:
    length: int
    type_code: int
    flags: int
    stream_id: int
    payload: bytes

    @property
    def type_name(self) -> str:
        return FRAME_TYPES.get(self.type_code, f"UNKNOWN(0x{self.type_code:02x})")

    @property
    def flag_names(self) -> str:
        known = FLAG_NAMES.get(self.type_code, {})
        names = [v for bit, v in known.items() if self.flags & bit]
        return "|".join(names) if names else "-"


def parse_frame(buf: bytes, offset: int) -> tuple[Frame, int]:
    """Parse one frame starting at offset. Raises on truncation."""
    if len(buf) - offset < 9:
        raise ValueError("truncated frame header")
    length = int.from_bytes(buf[offset : offset + 3], "big")
    type_code, flags = buf[offset + 3], buf[offset + 4]
    raw_sid = struct.unpack_from(">I", buf, offset + 5)[0]
    stream_id = raw_sid & 0x7FFFFFFF  # strip reserved bit
    end = offset + 9 + length
    if length > MAX_FRAME_SIZE:
        raise ValueError(f"frame length {length} exceeds MAX_FRAME_SIZE")
    if len(buf) < end:
        raise ValueError("truncated frame payload")
    return Frame(length, type_code, flags, stream_id,
                 buf[offset + 9 : end]), end


def build_fixture() -> bytes:
    """Preface + SETTINGS + SETTINGS-ACK + HEADERS + DATA on stream 1."""
    def frame(t: int, f: int, sid: int, payload: bytes) -> bytes:
        return (len(payload).to_bytes(3, "big") + bytes([t, f])
                + (sid & 0x7FFFFFFF).to_bytes(4, "big") + payload)

    settings = struct.pack(">HI", 0x3, 100)  # MAX_CONCURRENT_STREAMS=100
    headers_block = bytes([0x82, 0x86, 0x84])  # HPACK: :method GET,
                                               # :scheme http, :path /
    data = b'{"robot":"Weld-Bot"}'
    return b"".join([
        CLIENT_PREFACE,
        frame(0x4, 0x0, 0, settings),        # SETTINGS
        frame(0x4, 0x1, 0, b""),             # SETTINGS ack
        frame(0x1, 0x4 | 0x1, 1, headers_block),  # HEADERS END_HEADERS|END_STREAM
        frame(0x0, 0x1, 1, data),            # DATA END_STREAM
    ])


def main() -> None:
    stream = build_fixture()
    assert stream[:24] == CLIENT_PREFACE
    print(f"client preface (24 octets): {CLIENT_PREFACE!r}")
    print(f"{'len':>5} {'type':<14} {'flags':<28} {'stream':>6}")
    offset = len(CLIENT_PREFACE)
    frames: list[Frame] = []
    while offset < len(stream):
        frame_obj, offset = parse_frame(stream, offset)
        frames.append(frame_obj)
        print(f"{frame_obj.length:>5} {frame_obj.type_name:<14} "
              f"{frame_obj.flag_names:<28} {frame_obj.stream_id:>6}")
    assert [f.type_code for f in frames] == [0x4, 0x4, 0x1, 0x0]
    assert frames[2].stream_id == 1 and frames[2].flags & 0x4
    assert frames[3].payload == b'{"robot":"Weld-Bot"}'
    print("=== OK: 4 frames parsed, flags and stream IDs verified ===")


if __name__ == "__main__":
    main()
\end{minted}

\section{Common Pitfalls \& Anti-Patterns}
\label{sec:pitfalls}

\subsection*{P1 --- Content-Length / Transfer-Encoding Ambiguity (Request Smuggling)}

\begin{warningbox}
If a request carries \emph{both} \texttt{Content-Length} and \texttt{Transfer-Encoding: chunked}, RFC~9112~\S6.3 says the message is malformed and \texttt{Transfer-Encoding} overrides. When a front-end proxy honors \texttt{Content-Length} and a back-end honors \texttt{Transfer-Encoding} (or vice versa), an attacker can desynchronize the two parsers and \emph{smuggle} a second request into the connection --- the classic \textbf{CL.TE / TE.CL request smuggling} attack. Mitigations: reject ambiguous requests outright (\texttt{400}), normalize at the front proxy, and prefer HTTP/2 end-to-end (binary framing removes the ambiguity class entirely).
\end{warningbox}

\subsection*{P2 --- Assuming GET Never Has a Body}

RFC~9110 does not forbid a body on GET; it declares that a GET body has \emph{no generally defined semantics}. Some stacks (Elasticsearch historically) used GET bodies; many proxies drop or reject them. \textbf{Rule:} never put semantics-bearing content in a GET/HEAD/DELETE body; use POST if you need a request body on a read-like operation.

\subsection*{P3 --- Mishandling 100-continue}

A client sending \texttt{Expect: 100-continue} must be prepared for: (a)~an interim \texttt{100} followed by the final response, (b)~a final error (e.g.\ \texttt{413}) \emph{without} sending the body, or (c)~a server that never sends \texttt{100} --- in which case the client should transmit the body after a timeout (curl waits 1\,s). Clients that block forever waiting for \texttt{100}, or that treat the interim \texttt{100} as the final status, are both broken.

\subsection*{P4 --- Keep-Alive Timeout Mismatch}

If the client's idle-connection reuse window is \emph{longer} than the server/LB idle timeout, the client will race a stale connection: it sends a request onto a socket the middlebox just closed, producing intermittent \texttt{ECONNRESET}/``connection reset by peer'' errors under load. \textbf{Rule of thumb:} client idle timeout $<$ LB idle timeout $<$ server keep-alive timeout. This is the single most common cause of ``random'' production connection failures.

\subsection*{P5 --- Blindly Retrying Non-Idempotent Methods}

A timeout is ambiguous: the request may have been processed. Retrying a \texttt{POST /orders} can double-charge a customer. Only safe/idempotent methods (Table~\ref{tab:methods}) may be retried transparently; everything else needs application-level idempotency keys (Chapter~10) before any retry layer is switched on.

\subsection*{P6 --- Reading Until Close on a Keep-Alive Connection}

A client that reads an HTTP/1.1 response ``until EOF'' without honoring \texttt{Content-Length} or chunked framing hangs forever when the server keeps the connection open. Always parse framing headers first; read exactly the framed length.

\subsection*{P7 --- Ignoring Retry-After and Treating 5xx as Identical}

\texttt{503} + \texttt{Retry-After} is an explicit back-pressure signal; hammering it converts overload into outage. \texttt{429} without honoring \texttt{Retry-After} gets clients blacklisted. Distinguish transient (503, 429) from permanent (500, 501) and respect the server's stated delay with jittered backoff.

\subsection*{P8 --- Forgetting Hop-by-Hop Headers}

\texttt{Connection}, \texttt{Keep-Alive}, \texttt{TE}, \texttt{Trailer}, \texttt{Transfer-Encoding}, \texttt{Upgrade} are \emph{hop-by-hop}: proxies must not forward them. A hand-rolled reverse proxy that copies all headers verbatim leaks connection state across hops and breaks framing.

\section{Hands-On Production Lab \& Real-World Case Study}
\label{sec:lab}

\subsection{Lab A --- Dissect Real Traffic with curl.exe}

All commands run in PowerShell on Windows; use \texttt{curl.exe} explicitly so PowerShell does not alias \texttt{curl} to \texttt{Invoke-WebRequest}.

\begin{enumerate}
  \item \textbf{Start a local server.} In one terminal: \texttt{python -m http.server 8765 --bind 127.0.0.1}. It serves the current directory over plain HTTP/1.1.
  \item \textbf{Capture the verbose exchange.} In a second terminal:
\end{enumerate}

\begin{minted}{bash}
curl.exe -v http://127.0.0.1:8765/README.md -o NUL
\end{minted}

Lines prefixed \texttt{>} are the exact request bytes; \texttt{<} the response. Confirm: the request-line, the mandatory \texttt{Host} header, the \texttt{HTTP/1.1 200 OK} status-line, \texttt{Content-Length}, and the blank line separating head from body.

\begin{enumerate}[resume]
  \item \textbf{Force protocol versions and observe.}
\end{enumerate}

\begin{minted}{bash}
curl.exe -v --http1.1 http://127.0.0.1:8765/ -o NUL
curl.exe -v -I http://127.0.0.1:8765/README.md
\end{minted}

The \texttt{-I} issues a HEAD: identical headers, zero body bytes --- verify \texttt{Content-Length} is still present but no body arrives.

\begin{enumerate}[resume]
  \item \textbf{Conditional request.} Note the \texttt{ETag} (or \texttt{Last-Modified}) from step~2, then:
\end{enumerate}

\begin{minted}{bash}
curl.exe -v -H "If-None-Match: \"<etag>\"" http://127.0.0.1:8765/README.md -o NUL
\end{minted}

(The stdlib \texttt{SimpleHTTPRequestHandler} does not implement validators --- observe that it ignores the header and returns 200; repeat against any conditional-capable server, or your Listing-1 server extended with an \texttt{ETag}, to see \texttt{304 Not Modified}.)

\subsection{Lab B --- Your Raw Client vs.\ the Stdlib Server}

Run Listing~1, then re-run it with the request modified to omit \texttt{Connection: close} and observe that the client now hangs unless it honors \texttt{Content-Length} --- Pitfall~P6 reproduced on demand. Then run Listing~2 and match each printed chunk against the byte layout diagram in this chapter.

\subsection{Lab C --- Mint a Self-Signed Certificate and Inspect TLS}

\begin{minted}{bash}
# openssl.exe ships with Git for Windows:
# C:\Program Files\Git\usr\bin\openssl.exe
openssl.exe req -x509 -newkey rsa:2048 -keyout northwind_lab.key `
  -out northwind_lab.pem -days 365 -nodes `
  -subj "/CN=localhost/O=Northwind Robotics"
\end{minted}

Place \texttt{northwind\_lab.pem} beside Listing~3 and run it: the fixture decoder prints subject, issuer, SANs, and validity; the context audit prints the TLS~1.3 cipher list and ALPN set. (Self-signed certificates are for local labs only --- never disable verification in production to make one work.)

\subsection{Lab D --- HTTP/2 Frame Dissection}

Run Listing~4 and confirm the parsed sequence: \texttt{SETTINGS}, \texttt{SETTINGS(ACK)}, \texttt{HEADERS(END\_STREAM|END\_HEADERS)}, \texttt{DATA(END\_STREAM)} on stream~1. Then mutate the fixture: flip a flag bit, truncate a payload, and watch the parser's \texttt{ValueError} fire --- you now own a differential fuzzing harness.

\subsection{Real-World Case Study: The 60-Second Mystery}

\emph{War story (composite of common industry incidents, Northwind Robotics setting).} Northwind Robotics' telemetry API sat behind a cloud load balancer with a \textbf{60-second idle timeout}; the Python fleet's HTTP client pooled connections with a \textbf{300-second} keep-alive. Under bursty load, clients routinely grabbed connections the LB had silently reaped mid-idle. The first byte of the next request disappeared into a half-open socket; the client saw \texttt{ConnectionResetError} on $\sim$0.3\% of requests --- always the first request after an idle gap. Retries masked most failures, but a non-idempotent \texttt{POST /jobs} retry caused duplicate job dispatch (Pitfall~P5 compounding P4). Diagnosis took one packet capture: the LB's \texttt{FIN} arriving 60\,s after the last request, followed seconds later by the client's next request on the dead connection and an immediate \texttt{RST}. \textbf{Fix:} set the client pool's idle timeout to 45\,s (below the LB's 60\,s), raise the LB timeout to 75\,s, and gate \texttt{POST} retries behind idempotency keys. Errors dropped to zero; p99 latency fell 18\,ms because connection reuse actually worked instead of gambling.

\section{Self-Assessment Exercises \& Deep-Dive Questions}

\begin{enumerate}[label=\textbf{E1.\arabic*}, leftmargin=3.2em]
  \item \textbf{Byte archaeology.} Extend Listing~1 to send a POST with a \texttt{Content-Length} body, then deliberately declare a \texttt{Content-Length} one byte too small. Observe how the leftover byte corrupts the next exchange on a keep-alive connection. Explain the desynchronization in terms of framing.
  \item \textbf{Chunked boundaries.} Modify Listing~2's server to emit a chunked body where a chunk straddles two \texttt{recv} calls. Verify the client's \texttt{pending} buffer logic handles it; then reduce \texttt{MAX\_CHUNK} below a chunk's size and confirm the defensive \texttt{ValueError}.
  \item \textbf{Status-code audit.} For each endpoint of a fictional Northwind Robotics API (\texttt{POST /robots}, \texttt{GET /robots/\{id\}}, \texttt{PUT /robots/\{id\}/firmware}, \texttt{DELETE /robots/\{id\}}), choose the full set of success and error status codes you would return, with one-line justifications referencing Table~\ref{tab:statusdilemmas}.
  \item \textbf{Frame arithmetic.} A HEADERS frame carries a 70{,}000-octet header block. With the default \texttt{SETTINGS\_MAX\_FRAME\_SIZE}, how many frames are emitted, of which types, and with which flags set on the final frame?
  \item \textbf{HPACK reasoning.} A client sends 100 requests to the same API with identical headers except a per-request \texttt{x-trace-id}. Estimate the header bytes sent under HTTP/1.1 versus HTTP/2+HPACK, stating your assumptions about dynamic-table hits.
  \item \textbf{Flow-control deadlock hunt.} A client sets \texttt{SETTINGS\_INITIAL\_WINDOW\_SIZE = 0} mid-connection. Describe precisely which frames can still flow and why this is legal; then explain how a buggy server that never sends \texttt{WINDOW\_UPDATE} manifests to users.
  \item \textbf{TLS transcript.} From memory, list the TLS~1.3 handshake messages in order, mark which are encrypted, and identify the exact message after which the client could theoretically begin sending 0-RTT data on a resumption.
  \item \textbf{Smuggling sketch.} Draw the byte-level layout of a CL.TE smuggling payload against a front-end that trusts \texttt{Content-Length} and a back-end that trusts \texttt{Transfer-Encoding}. Then state the two mitigations that would have prevented it.
  \item \textbf{HOL cost model.} Given 1\% uniform packet loss, 50\,ms RTT, and 10 parallel small requests, argue qualitatively when HTTP/2 over one connection loses to HTTP/1.1 over six, and why HTTP/3's outcome differs.
  \item \textbf{Deep dive --- 0-RTT danger.} Explain why 0-RTT early data is replayable, why that makes \texttt{POST /payments} over 0-RTT unacceptable, and which defences (single-use tickets, early-data allowlists of safe methods) are appropriate.
\end{enumerate}

\section{Interview Q\&A Vault}

\subsection*{Junior Level}

\begin{enumerate}[label=\textbf{Q\arabic*.}, leftmargin=2.8em]
  \item \textbf{What are the three parts of an HTTP/1.1 request on the wire?} The request-line (\texttt{METHOD SP target SP HTTP/1.1}), zero or more header fields each ending in CRLF, a mandatory empty line (CRLF), and an optional body delimited by \texttt{Content-Length} or chunked framing.
  \item \textbf{What bytes terminate every HTTP/1.1 header line?} CRLF --- carriage return + line feed, \texttt{0x0D 0x0A}. Not LF alone; bare-LF tolerance is a parser quirk that smuggling attacks exploit.
  \item \textbf{Why is the \texttt{Host} header mandatory in HTTP/1.1?} It enables virtual hosting: one IP serves many domains, and the server (and TLS SNI routing) needs the name to pick the right site. Absence $\Rightarrow$ \texttt{400 Bad Request}.
  \item \textbf{GET vs.\ POST in one sentence?} GET retrieves a representation and is safe, idempotent, and cacheable; POST asks the resource to process the enclosed body and has none of those guarantees.
  \item \textbf{What does idempotent mean, precisely?} Repeating the identical request $n$ times has the same intended effect on server state as one request. It says nothing about identical \emph{responses} (a retried DELETE may return 404 the second time and still be idempotent).
  \item \textbf{What is the difference between 401 and 403?} 401 = unauthenticated: credentials missing or invalid, and the response must carry \texttt{WWW-Authenticate}. 403 = authenticated (or authentication irrelevant) but not allowed; re-authenticating will not help.
  \item \textbf{When do you return 201 instead of 200?} When the request created one or more new resources; include a \texttt{Location} header pointing at the primary new resource. 200 is for generic success with a payload.
  \item \textbf{What is 204 used for?} Success with no representation to return (e.g.\ a DELETE or an idempotent save). A 204 response must not carry a body.
  \item \textbf{What does \texttt{Content-Type} declare vs.\ \texttt{Accept}?} \texttt{Content-Type} describes the body actually present in this message; \texttt{Accept} is a request header describing what the client can parse in the response.
  \item \textbf{What is a q-value?} A relative weight in $[0,1]$ on an \texttt{Accept} media range, e.g.\ \texttt{application/json;q=0.9}. The server serves the highest-$q$ representation it can produce, else \texttt{406}.
  \item \textbf{What is keep-alive?} HTTP/1.1's default persistent-connection behavior: the TCP connection stays open after the response so subsequent requests skip connection setup. It is why message framing (\texttt{Content-Length}/chunked) matters.
  \item \textbf{What does \texttt{Connection: close} do?} Tells the peer to close the TCP connection after this message; for responses it also makes ``read until EOF'' a legal body framing.
\end{enumerate}

\subsection*{Mid Level}

\begin{enumerate}[label=\textbf{Q\arabic*.}, leftmargin=2.8em, start=13]
  \item \textbf{Why did HTTP/1.1 pipelining die?} Responses must arrive in request order, so one slow response blocks the whole pipeline (HOL blocking); intermediaries corrupted pipelined traffic; and a failed pipeline cannot safely retry non-idempotent requests. Browsers shipped it disabled and removed it.
  \item \textbf{Explain chunked transfer encoding byte layout.} Each chunk: ASCII-hex size, CRLF, that many data octets, CRLF. Terminated by a zero-size chunk, optional trailers, final CRLF. It enables streaming bodies of unknown length over persistent connections.
  \item \textbf{400 vs.\ 422?} 400: the server could not parse the request --- malformed syntax, bad framing. 422: the request parsed fine but its semantics fail (e.g.\ JSON schema violations, a negative quantity). 400 blames the bytes; 422 blames the meaning.
  \item \textbf{409 vs.\ 412?} 409: the request conflicts with current resource state (duplicate key, version mismatch on an unconditional write). 412: a conditional request's precondition header (\texttt{If-Match}, \texttt{If-Unmodified-Since}) evaluated false --- the client asked for precondition checking and lost.
  \item \textbf{Why send \texttt{Retry-After} with 429 and 503?} It converts client retry behavior from a thundering herd into a scheduled drain: the server tells clients the earliest sensible retry moment. Without it, retry storms deepen the overload that caused the 503.
  \item \textbf{How does \texttt{If-None-Match} save bandwidth?} The client echoes a stored \texttt{ETag}; if the representation is unchanged the server returns \texttt{304 Not Modified} with no body --- headers only. On a safe method it is a cache revalidation; with \texttt{If-Match} on PUT it becomes optimistic concurrency control.
  \item \textbf{What does \texttt{Cache-Control: no-cache} actually mean?} ``Store is allowed, but revalidate with the origin before every reuse.'' It is chronically confused with \texttt{no-store} (never persist at all). Quoting the wrong one either leaks sensitive data or destroys cache efficiency.
  \item \textbf{What is CORS preflight and who enforces it?} For non-simple cross-origin requests the browser sends \texttt{OPTIONS} with \texttt{Origin} + \texttt{Access-Control-Request-Method}; the server's \texttt{Access-Control-Allow-*} headers authorize the real request. The \emph{browser} enforces it --- CORS is not a server-side security boundary; curl ignores it.
  \item \textbf{Describe the HTTP/2 frame header.} Nine octets: 24-bit payload length, 8-bit type, 8-bit flags, 1 reserved bit + 31-bit stream identifier. Everything else in HTTP/2 is payload typed by those fields.
  \item \textbf{Name all ten HTTP/2 frame types and their codes.} DATA=\texttt{0x0}, HEADERS=\texttt{0x1}, PRIORITY=\texttt{0x2}, RST\_STREAM=\texttt{0x3}, SETTINGS=\texttt{0x4}, PUSH\_PROMISE=\texttt{0x5}, PING=\texttt{0x6}, GOAWAY=\texttt{0x7}, WINDOW\_UPDATE=\texttt{0x8}, CONTINUATION=\texttt{0x9}.
  \item \textbf{What is the HTTP/2 connection preface?} The client magic \texttt{PRI * HTTP/2.0\textbackslash r\textbackslash n\textbackslash r\textbackslash nSM\textbackslash r\textbackslash n\textbackslash r\textbackslash n} (24 octets) followed by a SETTINGS frame; it lets a server distinguish h2c from HTTP/1.1 on a shared port and seeds connection parameters.
  \item \textbf{How does HTTP/2 multiplexing work?} One TCP connection carries many streams; every frame is tagged with a stream ID, so requests/responses interleave at frame granularity and one slow response no longer blocks others. Concurrency limits come from \texttt{SETTINGS\_MAX\_CONCURRENT\_STREAMS}.
  \item \textbf{What is HPACK and why does HTTP/2 need it?} HPACK compresses header blocks using a 61-entry static table, a per-connection dynamic table of recently used headers (default 4096 octets), and static Huffman coding. Without it, per-request textual headers (cookies alone can be kilobytes) would dominate small-payload traffic and defeat TCP's initial congestion window.
  \item \textbf{Why must HPACK header blocks be decoded in order, and how does QPACK fix it?} HPACK's dynamic table is a shared ordered state machine; decoding block $n$ requires having applied blocks $<n$. On QUIC's independent streams that ordering would re-couple streams, so QPACK moves table updates onto separate encoder streams and tracks acknowledgements, tolerating out-of-order delivery at the cost of possible encoder-delay stalls.
\end{enumerate}

\subsection*{Senior Level}

\begin{enumerate}[label=\textbf{Q\arabic*.}, leftmargin=2.8em, start=27]
  \item \textbf{Explain HTTP/2 flow control.} Credit-based, hop-by-hop: every DATA octet consumes credit from both the stream's window and the connection window (both default 65{,}535). Receivers replenish with \texttt{WINDOW\_UPDATE} increments ($1..2^{31}-1$); senders resize peer stream windows via \texttt{SETTINGS\_INITIAL\_WINDOW\_SIZE}. Only DATA counts; a zero window stalls DATA but never HEADERS/PING/RST, so deadlocks are recoverable.
  \item \textbf{What HOL blocking remains in HTTP/2?} TCP's ordered byte stream: one lost segment stalls all streams until retransmission, because TCP cannot deliver later bytes early. Multiplexing removed application-layer HOL only; the transport layer re-imposed it. On lossy links this can make h2 slower than parallel h1 connections.
  \item \textbf{Why does HTTP/3 use UDP --- isn't UDP unreliable?} UDP is deliberately minimal: QUIC reimplements reliability, ordering, and congestion control \emph{in userspace} on top of it. Kernel TCP and middlebox ossification make deploying a new TCP-like transport impossible; UDP is the only ubiquitous substrate that permits transport evolution, per-stream independence, and fused TLS~1.3.
  \item \textbf{How does QUIC connection migration work?} Connections are bound to Connection IDs, not the 5-tuple. When a client's IP/port changes (Wi-Fi$\to$LTE), packets with the same Connection ID still identify the connection; the peer validates the new path (path validation challenge) and continues --- no re-handshake, streams resume where they left off.
  \item \textbf{TLS 1.3 handshake: which messages exist and which are encrypted?} \texttt{ClientHello} and \texttt{ServerHello} are cleartext (plus legacy record wrapping); from \texttt{EncryptedExtensions} onward --- \texttt{Certificate}, \texttt{CertificateVerify}, \texttt{Finished} --- everything is encrypted under handshake traffic keys derived from the ECDHE key share exchange completed by ServerHello.
  \item \textbf{TLS 1.2 vs.\ 1.3 round trips?} 1.2 full handshake: 2-RTT; 1.3: 1-RTT because the client speculatively sends its \texttt{key\_share} in ClientHello. Resumption: 1.2 gives 1-RTT; 1.3 gives 0-RTT early data --- which is replayable and therefore must be restricted to idempotent requests.
  \item \textbf{What do SNI and ALPN each negotiate?} SNI (\texttt{server\_name}) selects which certificate identity the server presents on a shared IP. ALPN selects the application protocol (\texttt{h2}, \texttt{http/1.1}) inside the handshake, avoiding a post-handshake negotiation round trip; failure to agree aborts with \texttt{no\_application\_protocol}.
  \item \textbf{Walk certificate chain validation.} Verify the leaf's hostname against SANs, validity window, and signature using the presented intermediate(s); recurse until a locally trusted root anchor; enforce \texttt{basicConstraints CA=TRUE} and key-usage on intermediates; optionally check revocation (OCSP/CRL). Any failure aborts before the first request byte.
  \item \textbf{What is request smuggling and how does HTTP/2 help?} Disagreement between two HTTP/1.1 parsers about message boundaries (CL.TE/TE.CL desync) lets an attacker prefix a hidden request onto the next victim's request. HTTP/2's binary length-prefixed framing removes textual ambiguity --- though h2$\to$h1 downgrades reintroduce risk, so gateways must normalize strictly.
  \item \textbf{Why was HTTP/2 server push deprecated?} Adoption was negligible and often harmful: servers pushed assets clients already cached, wasting bandwidth; cache-aware push never materialized; priority interactions were buggy. RFC~9113 deprecates it, browsers removed support, and \texttt{103 Early Hints} covers the honest use case.
  \item \textbf{How would you debug intermittent connection-reset errors behind an LB?} Suspect a keep-alive race first: compare client pool idle timeout, LB idle timeout, and server keep-alive timeout. Confirm with a packet capture showing \texttt{FIN} from the LB followed by a client request and \texttt{RST}. Fix by making the client timeout strictly shorter than the LB's.
  \item \textbf{Design the retry policy for \texttt{POST /payments} over flaky HTTP/2.} Never transport-retry blindly. Require an idempotency key, retry only on connection-refused/pre-send failures without it, bound retries, honor \texttt{Retry-After}/429, jitter backoff, and make the server deduplicate on the key. The protocol tells you POST is not idempotent; the application must restore the property.
  \item \textbf{When does 0-RTT help and when does it hurt?} It eliminates handshake latency for resumed connections --- valuable for mobile/high-RTT clients. It hurts when early data is replayed by a network attacker: restrict 0-RTT to safe, idempotent requests, use single-use tickets, and keep anti-replay windows tight.
  \item \textbf{Why do proxies strip \texttt{Connection} and \texttt{Keep-Alive}?} They are hop-by-hop headers describing one link, not the path. Forwarding them would assert connection properties across hops the header knows nothing about, corrupting framing and tunnel semantics (RFC~9110~\S7.6.1; HTTP/2 forbids them outright).
\end{enumerate}

\section{External Bibliography \& Further Reading}

\subsection*{Normative Specifications (read these first)}
\begin{itemize}
  \item RFC~9110, \emph{HTTP Semantics} \cite{rfc9110} --- methods, status codes, header fields, content negotiation, conditional requests, caching semantics.
  \item RFC~9112, \emph{HTTP/1.1} \cite{rfc9112} --- message syntax, framing, chunked coding, connection management.
  \item RFC~9113, \emph{HTTP/2} \cite{rfc9113} --- binary framing, streams, flow control, push deprecation.
  \item RFC~9114, \emph{HTTP/3} \cite{rfc9114} --- HTTP semantics over QUIC streams.
  \item RFC~9000, \emph{QUIC} \cite{rfc9000} --- the transport: packets, streams, loss recovery, connection IDs, migration.
  \item RFC~8446, \emph{TLS 1.3} \cite{rfc8446} --- handshake messages, key schedule, 0-RTT.
  \item RFC~7541, \emph{HPACK} \cite{rfc7541} --- header compression for HTTP/2.
  \item RFC~9204, \emph{QPACK} (inline; not in the shared bibliography) --- field compression for HTTP/3.
  \item RFC~7231 and RFC~7540 (superseded by 9110/9113 but ubiquitous in older literature; inline reference).
\end{itemize}

\subsection*{Books}
\begin{itemize}
  \item Barry Pollard, \emph{HTTP/2 in Action}, Manning, 2019 (inline reference) --- the most practical h2 operations book; covers framing, HPACK, and real deployment pitfalls.
  \item Ilya Grigorik, \emph{High Performance Browser Networking}, O'Reilly, 2013 (free at \texttt{hpbn.co}; inline reference) --- TCP/TLS/HTTP latency mechanics; the canonical treatment of why round trips dominate.
  \item Martin Kleppmann, \emph{Designing Data-Intensive Applications} \cite{kleppmann2017ddia} --- Chapter~1 and the distributed-systems framing for retries, idempotence, and exactly-once illusions.
  \item Michael T.\ Nygard, \emph{Release It!}, 2nd ed.\ \cite{nygard2018releaseit} --- timeout, connection-pool, and stability patterns; the source of the keep-alive-mismatch failure genus.
  \item Roy T.\ Fielding, \emph{Architectural Styles and the Design of Network-based Software Architectures} \cite{fielding2000rest} --- the dissertation; Chapter~5 defines REST, Chapter~6 motivates HTTP's design.
\end{itemize}

\subsection*{Documentation, Tools \& Courses}
\begin{itemize}
  \item Mozilla MDN Web Docs, \emph{HTTP} reference (inline) --- best-maintained header/status-code reference.
  \item Wireshark TLS/QUIC dissectors and the \texttt{sslkeylog} (\texttt{SSLKEYLOGFILE}) workflow (inline) --- decrypt your own lab traffic.
  \item \texttt{nghttp2} (\texttt{nghttp}, \texttt{h2load}) and \texttt{curl --http2/--http3} documentation (inline) --- reference clients for protocol-version experiments.
  \item Cloudflare Learning Center and the \emph{Smashing Magazine}/Cloudflare QUIC deployment post-mortems (inline) --- real-world HTTP/3 rollout data.
\end{itemize}

\section{Capstone Projects}

\subsection{Capstone 1 --- From-Scratch HTTP/1.1 Server [junior]}
\textbf{Objective:} Implement a minimal but correct HTTP/1.1 server over raw \texttt{socketserver}, serving the fictional Northwind Robotics parts catalog.\\
\textbf{Inputs/Datasets:} an in-repo JSON file \texttt{parts.json} (synthetic catalog, $\leq$50 entries).\\
\textbf{Steps:} (1)~parse request-line and headers until CRLFCRLF; (2)~implement GET/HEAD for \texttt{/parts} and \texttt{/parts/\{id\}}; (3)~emit correct \texttt{Content-Length}, \texttt{Content-Type}, \texttt{Connection} semantics; (4)~support keep-alive for at least 5 sequential requests per connection; (5)~return 404/400/405 correctly.\\
\textbf{Acceptance Criteria:}
\begin{itemize}
  \item[$\square$] \texttt{curl.exe -v} shows byte-correct status-line, headers, and framing.
  \item[$\square$] Ten sequential keep-alive requests on one connection all succeed.
  \item[$\square$] Malformed request-line yields 400; unknown path yields 404; POST yields 405.
  \item[$\square$] HEAD returns headers identical to GET with zero body bytes.
\end{itemize}
\textbf{Stretch Goals:} conditional GET with \texttt{ETag}/\texttt{If-None-Match} $\to$ 304; graceful \texttt{Connection: close}.

\subsection{Capstone 2 --- Chunked Streaming Proxy [mid]}
\textbf{Objective:} Build a localhost reverse proxy that relays responses using \texttt{Transfer-Encoding: chunked} while enforcing defensive parsing.\\
\textbf{Inputs/Datasets:} Listing-2-style local origin server emitting chunked telemetry; \texttt{telemetry.csv} synthetic dataset.\\
\textbf{Steps:} (1)~accept client connections; (2)~forward request head verbatim; (3)~re-chunk origin bytes in 256-octet chunks; (4)~enforce max chunk bounds and CRLF validation; (5)~log per-request byte counts with \texttt{logging}.\\
\textbf{Acceptance Criteria:}
\begin{itemize}
  \item[$\square$] Client receives identical decoded body as direct origin fetch.
  \item[$\square$] Oversized or malformed chunk from origin is rejected with 502 and logged.
  \item[$\square$] Hop-by-hop headers are stripped per RFC~9110~\S7.6.1.
  \item[$\square$] Runs fully offline; deterministic test script included.
\end{itemize}
\textbf{Stretch Goals:} trailer support (\texttt{Trailer: X-Checksum}); backpressure via bounded buffers.

\subsection{Capstone 3 --- HPACK-Lite Header Encoder [mid]}
\textbf{Objective:} Implement a pure-stdlib encoder for a subset of HPACK: static-table indexing, literal-with-incremental-indexing, and integer prefix coding (RFC~7541 \S5).\\
\textbf{Inputs/Datasets:} a fixed request-header corpus (\texttt{headers\_corpus.json}, synthetic Northwind Robotics API calls).\\
\textbf{Steps:} (1)~implement RFC~7541 integer encoding ($N$-bit prefixes); (2)~match against the 61-entry static table; (3)~maintain a bounded dynamic table with eviction; (4)~emit hex dumps per header block; (5)~measure bytes saved vs.\ HTTP/1.1 text across the corpus.\\
\textbf{Acceptance Criteria:}
\begin{itemize}
  \item[$\square$] Integer encoder round-trips all corpus values, including prefix-overflow cases.
  \item[$\square$] Second occurrence of an identical header block encodes as a single indexed reference.
  \item[$\square$] Dynamic table never exceeds its configured size bound.
  \item[$\square$] Report shows $\geq$60\% byte reduction on the repeat-heavy corpus.
\end{itemize}
\textbf{Stretch Goals:} static Huffman coding; a decoder that round-trips your encoder.

\subsection{Capstone 4 --- Protocol Fingerprinting \& Frame Dissection Toolkit [senior]}
\textbf{Objective:} Extend Listing~4 into a passive analysis CLI that ingests captured byte streams (files), classifies the protocol (HTTP/1.1 text, h2 preface, TLS records), and emits a structured report.\\
\textbf{Inputs/Datasets:} a \texttt{captures\textbackslash} folder of synthetic fixtures generated by your own scripts (HTTP/1.1 exchanges, h2 frame sequences, TLS record headers).\\
\textbf{Steps:} (1)~sniff protocol by magic bytes/preface; (2)~dispatch to per-protocol parsers; (3)~for h2, emit per-frame CSV (type, flags, stream, length); (4)~detect anomalies: CL+TE ambiguity, frames exceeding 16\,KiB, RST floods; (5)~exit codes suitable for CI.\\
\textbf{Acceptance Criteria:}
\begin{itemize}
  \item[$\square$] Correctly classifies all provided fixtures with zero false positives.
  \item[$\square$] Flags a planted CL.TE smuggling capture and a truncated frame.
  \item[$\square$] Output is deterministic (stable ordering, no timestamps) for diff-based tests.
  \item[$\square$] \texttt{--json} mode emits machine-readable reports.
\end{itemize}
\textbf{Stretch Goals:} HPACK state reconstruction across HEADERS/CONTINUATION sequences; entropy-based flags for non-HTTP payloads.

\subsection{Capstone 5 --- HTTP Version Latency Profiler \& Wire-Debugging Workbench [staff]}
\textbf{Objective:} Build a local benchmark harness that quantifies connection-setup and per-request costs across HTTP/1.1 configurations (cold connection, keep-alive warm, pipelining simulation, chunked vs.\ fixed-length), producing the standard addendum's benchmark table from real measurements.\\
\textbf{Inputs/Datasets:} your own local servers from Capstones 1--2; synthetic payload set (1\,KB, 64\,KB, 1\,MB) in \texttt{payloads\textbackslash}.\\
\textbf{Steps:} (1)~harness spins up servers on ephemeral ports; (2)~measure p50/p99 for each configuration with seeded, fixed iteration counts (no wall-clock assertions in tests --- record, don't assert); (3)~simulate loss/latency with a socket-level delay shim; (4)~emit the benchmark table as Markdown + CSV; (5)~document how the methodology would extend to h2/h3 with the \texttt{h2}/\texttt{aioquic} libraries.\\
\textbf{Acceptance Criteria:}
\begin{itemize}
  \item[$\square$] Harness runs offline end-to-end with one command and seeds all randomness.
  \item[$\square$] Report quantifies keep-alive benefit (warm vs.\ cold) for each payload size.
  \item[$\square$] Delay shim demonstrates HOL-blocking cost growth under injected latency.
  \item[$\square$] Benchmark table conforms to the chapter addendum template.
\end{itemize}
\textbf{Stretch Goals:} pluggable transport interface so a real h2 client (via \texttt{h2}) drops into the same harness; flame-graph-friendly structured logs.

\section{Python Dependencies Appendix}

The core listings in this chapter are \textbf{stdlib-only}. One optional dependency extends the laboratories:

\begin{minted}{text}
# requirements.txt  (chapter 1 - optional tooling only; core is stdlib)
h2>=4.1.0
\end{minted}

\subsection{h2 ($\geq$4.1.0)}
\textbf{Description:} A pure-Python, in-memory HTTP/2 protocol state machine (hyper-h2). It handles framing, HPACK, and flow control; you supply sockets and TLS.\\
\textbf{Why included:} Listings 1--4 teach the wire by hand; \texttt{h2} is the bridge to speaking \emph{real} h2 against local servers in the stretch goals and Capstone~5 without implementing the full protocol stack.\\
\textbf{Alternatives:} \texttt{httpx[http2]} (full client, hides the frames --- pedagogically backwards here); \texttt{aioquic} (if you extend to HTTP/3); raw \texttt{hyperframe} (lower level than needed).\\
\textbf{Installation (PowerShell):}
\begin{minted}{bash}
python -m venv .venv
.\.venv\Scripts\Activate.ps1
pip install -r requirements.txt
\end{minted}

\section{Chapter Summary \& Key Takeaways}

\begin{itemize}
  \item HTTP/1.1 is a textual protocol framed by CRLF; every message boundary derives from \texttt{Content-Length}, chunked coding, or connection close --- and ambiguity between the first two is a security vulnerability class (request smuggling), not a style question \cite{rfc9112}.
  \item Method properties (safe/idempotent/cacheable) and status-code precision (201 vs.\ 200 vs.\ 204; 400 vs.\ 422; 401/403/404; 409 vs.\ 412; 429/503 + \texttt{Retry-After}) are contracts that caches, retries, and intermediaries enforce --- get them right and infrastructure works \emph{for} you \cite{rfc9110}.
  \item HTTP/2 replaces text with a 9-octet frame header (24-bit length, 8-bit type, 8-bit flags, 31-bit stream ID), multiplexes streams over one TCP connection, compresses headers with HPACK, and flow-controls with dual windows --- but inherits TCP's transport-layer HOL blocking \cite{rfc9113,rfc7541}.
  \item HTTP/3/QUIC moves streams into userspace over UDP: per-stream loss independence, Connection-ID-based migration, and a fused 1-RTT TLS~1.3 handshake with 0-RTT resumption for repeat clients \cite{rfc9114,rfc9000,rfc8446}.
  \item TLS~1.3's six-message handshake, mandatory forward secrecy, and ALPN-negotiated protocol selection define both the security and the latency floor of every HTTPS call; keep-alive hit rate is therefore an economic metric, not a nicety.
  \item The dominant production failure modes are protocol misunderstandings: keep-alive timeout mismatches, blind non-idempotent retries, read-until-close framing bugs, and ignored back-pressure signals. Each is preventable with the mental models --- and the four runnable tools --- built in this chapter.
\end{itemize}

\section{Quick-Reference Card}

\begin{table}[h]
  \centering
  \small
  \begin{tabularx}{\textwidth}{l l X}
    \toprule
    \textbf{Method} & \textbf{Safe/Idem./Cache} & \textbf{Use for} \\
    \midrule
    GET     & Y / Y / Y & Read a resource. \\
    HEAD    & Y / Y / Y & Metadata only (no body). \\
    POST    & N / N / rare & Create / process; needs idempotency key to retry. \\
    PUT     & N / Y / N & Full replacement at known URI. \\
    PATCH   & N / N / N & Partial update. \\
    DELETE  & N / Y / N & Remove; 204 on success. \\
    OPTIONS & Y / Y / N & Capabilities / CORS preflight. \\
    TRACE   & Y / Y / N & Echo diagnostics (usually disabled). \\
    CONNECT & N / N / N & Tunnel establishment. \\
    \bottomrule
  \end{tabularx}
  \caption{Method cheat sheet.}
\end{table}

\begin{table}[h]
  \centering
  \small
  \begin{tabularx}{\textwidth}{l X}
    \toprule
    \textbf{Code} & \textbf{Meaning / when to emit} \\
    \midrule
    200 / 201 / 204 & OK with body / created (+Location) / OK, no body. \\
    301 / 302 / 307 / 308 & Permanent / temporary redirect; 307/308 preserve method+body. \\
    304 & Not modified (conditional GET hit). \\
    400 / 401 / 403 / 404 & Malformed / unauthenticated / forbidden / not found (or masked). \\
    405 / 406 / 409 / 412 / 413 / 415 / 422 / 429 & Method not allowed / not acceptable / conflict / precondition failed / payload too large / unsupported media / semantic error / rate limited (+Retry-After). \\
    500 / 502 / 503 / 504 & Server fault / bad gateway / unavailable (+Retry-After) / gateway timeout. \\
    \bottomrule
  \end{tabularx}
  \caption{Status-code cheat sheet.}
\end{table}

\begin{table}[h]
  \centering
  \small
  \begin{tabularx}{\textwidth}{l l X}
    \toprule
    \textbf{Code} & \textbf{Frame} & \textbf{One-liner} \\
    \midrule
    \texttt{0x0} & DATA & Body octets; counts against both flow-control windows. \\
    \texttt{0x1} & HEADERS & HPACK block; opens a stream. \\
    \texttt{0x2} & PRIORITY & Deprecated dependency/weight hint. \\
    \texttt{0x3} & RST\_STREAM & Kill one stream (32-bit error code). \\
    \texttt{0x4} & SETTINGS & Connection parameters; must be ACKed. \\
    \texttt{0x5} & PUSH\_PROMISE & Push reservation (deprecated in practice). \\
    \texttt{0x6} & PING & Liveness / RTT (8 octets, ACK echo). \\
    \texttt{0x7} & GOAWAY & Graceful shutdown; last-stream-id. \\
    \texttt{0x8} & WINDOW\_UPDATE & Flow-control credit increment. \\
    \texttt{0x9} & CONTINUATION & Header-block spillover; needs END\_HEADERS. \\
    \bottomrule
  \end{tabularx}
  \caption{HTTP/2 frame-type cheat sheet (header: 24-bit length, 8-bit type, 8-bit flags, 1 reserved bit, 31-bit stream ID).}
\end{table}

\section{Standardized Evidence Addendum}
\subsection{Benchmark Snapshot Template}
\begin{table}[h]
  \centering
  \small
  \begin{tabularx}{\textwidth}{l c c c c}
    \toprule
    \textbf{Config} & \textbf{p50 latency} & \textbf{p99 latency} & \textbf{Error rate} & \textbf{Throughput} \\
    \midrule
    Baseline & -- & -- & -- & 1.00x \\
    Candidate A & -- & -- & -- & -- \\
    Candidate B & -- & -- & -- & -- \\
    \bottomrule
  \end{tabularx}
  \caption{Minimum benchmark table to keep per chapter.}
\end{table}

\subsection{Request Trace Template}
\begin{itemize}
  \item \textbf{Request:} one representative API call for this chapter's topic.
  \item \textbf{Before:} behavior (headers, payload, latency, failure mode) with the naive approach.
  \item \textbf{After:} behavior with the technique in this chapter.
  \item \textbf{Result:} what changed in correctness, resilience, latency, and operability.
\end{itemize}

\subsection{Cost and Latency Budget Snapshot}
\begin{table}[h]
  \centering
  \small
  \begin{tabularx}{\textwidth}{l c c}
    \toprule
    \textbf{Stage} & \textbf{Latency budget} & \textbf{Cost driver} \\
    \midrule
    TLS / connection setup & -- & handshake round trips, keep-alive hit rate \\
    Request validation & -- & schema complexity, CPU per request \\
    Business logic and I/O & -- & downstream fan-out, query cost \\
    Serialization and egress & -- & payload size, compression ratio \\
    \bottomrule
  \end{tabularx}
  \caption{Operational budget template for this chapter's technique.}
\end{table}

\section{Configuration Preset Template}
\begin{minted}{text}
# api_policy.yaml
policy_version: "v1.0.0"
component: "<chapter_topic>"
timeout_ms: 0
retry_budget: 0
fallback_strategy: "fail_fast"
rollback_trigger: "error_budget_burn_gt_threshold"
\end{minted}

\section{CI Quality Gate Specification}
\begin{itemize}
  \item contract tests green against the pinned OpenAPI/AsyncAPI/Protobuf spec,
  \item no breaking schema changes without a version bump,
  \item error responses conform to RFC 9457 problem-details shape,
  \item benchmark deltas within approved regression bounds (latency and throughput),
  \item policy version and rollback plan present in release metadata.
\end{itemize}

\section{Incident Runbook Template}
\begin{enumerate}
  \item Detect regression from SLO burn-rate alerts or consumer reports.
  \item Scope impacted endpoints, versions, and consumer segments.
  \item Contain impact (rollback, feature flag, rate-limit tightening, or traffic shift).
  \item Diagnose with traces, structured logs, and contract diffs.
  \item Recover with validated fix and verified rollout procedure.
  \item Prevent recurrence via new regression tests and CI gates.
\end{enumerate}

\section{Cross-Chapter Decision Card}
\begin{table}[h]
  \centering
  \small
  \begin{tabularx}{\textwidth}{l X X}
    \toprule
    \textbf{Dimension} & \textbf{Choose this approach when} & \textbf{Prefer an alternative when} \\
    \midrule
    Use case & -- & -- \\
    Latency budget & -- & -- \\
    Operational cost & -- & -- \\
    Team maturity & -- & -- \\
    \bottomrule
  \end{tabularx}
  \caption{Decision card to complete per chapter.}
\end{table}

\section{ROI Model and Build-vs-Buy}
\begin{itemize}
  \item \textbf{Build when:} the capability is differentiating, compliance forbids vendors, or scale makes vendor pricing irrational.
  \item \textbf{Buy when:} the capability is commodity (auth, gateways, observability), time-to-market dominates, or staffing is thin.
  \item \textbf{Hidden costs to model:} on-call load, upgrade cadence, expertise ramp, data egress fees.
\end{itemize}

\section{Disaster Recovery Notes (RPO/RTO)}
\begin{itemize}
  \item Define recovery point objective (RPO): maximum acceptable data loss window.
  \item Define recovery time objective (RTO): maximum acceptable restoration time.
  \item Test restores regularly; an untested backup is not a backup.
\end{itemize}


